\begin{abstract}
Claims organisations must decide which documents to request from a claimant; one fact can make a document required, unnecessary or premature. Language-model agents usually apply such rules themselves, from their prompt. This article examines whether code should apply them instead, to the agent's verified judgments of the case. CasePath records a true, false or unresolved judgment with a quotation for every rule condition, verifies these judgments and executes source-grounded process rules that request, hold or skip each document. Its requests are provably never premature, and exact deferral removes an incompleteness of its three-valued process engine. CasePath was compared with five alternatives on 90 tenancy claims and 45 pairs of theft cases, using two current models and three runs per case. With the cost-efficient model configured in the product, execution made four times fewer unsafe errors on complete request lists and four times fewer wrong changes when a fact changed. With a frontier model, rules in the prompt were equally safe and more accurate on automatically extracted rules, whose errors execution reproduces. Executing rules makes request quality robust to the model at the price of more clarification questions; its accuracy is bounded by the quality of its rules.
\end{abstract}
